% Q13/Q16 -- 4:1 MUX using nMOS pass-transistor logic (schematic).
% 8 nMOS in a 4x2 array; two series transistors per row; outputs tied to f.
\documentclass[border=10pt]{standalone}
\usepackage{circuitikz}
\begin{document}
\begin{circuitikz}
% output bus
\draw (8,1.0) -- (8,7.0);
\draw (8,4.0) -- (9.2,4.0) node[right]{$f$};
% rows: I0..I3 at decreasing y; gate labels (g1 then g2)
% row I0 : Sbar1, Sbar0
\path (3,7) node[nmos, rotate=-90] (a0){};
\path (6,7) node[nmos, rotate=-90] (b0){};
\draw (a0.S) -- (1.2,7) node[left]{$I_0$};
\draw (a0.D) -- (b0.S);
\draw (b0.D) -- (8,7);
\draw (a0.G) -- (3,8.0) node[above]{$\overline{S_1}$};
\draw (b0.G) -- (6,8.0) node[above]{$\overline{S_0}$};
% row I1 : Sbar1, S0
\path (3,5) node[nmos, rotate=-90] (a1){};
\path (6,5) node[nmos, rotate=-90] (b1){};
\draw (a1.S) -- (1.2,5) node[left]{$I_1$};
\draw (a1.D) -- (b1.S);
\draw (b1.D) -- (8,5);
\draw (a1.G) -- (3,5.9) node[above]{$\overline{S_1}$};
\draw (b1.G) -- (6,5.9) node[above]{$S_0$};
% row I2 : S1, Sbar0
\path (3,3) node[nmos, rotate=-90] (a2){};
\path (6,3) node[nmos, rotate=-90] (b2){};
\draw (a2.S) -- (1.2,3) node[left]{$I_2$};
\draw (a2.D) -- (b2.S);
\draw (b2.D) -- (8,3);
\draw (a2.G) -- (3,3.9) node[above]{$S_1$};
\draw (b2.G) -- (6,3.9) node[above]{$\overline{S_0}$};
% row I3 : S1, S0
\path (3,1) node[nmos, rotate=-90] (a3){};
\path (6,1) node[nmos, rotate=-90] (b3){};
\draw (a3.S) -- (1.2,1) node[left]{$I_3$};
\draw (a3.D) -- (b3.S);
\draw (b3.D) -- (8,1);
\draw (a3.G) -- (3,1.9) node[above]{$S_1$};
\draw (b3.G) -- (6,1.9) node[above]{$S_0$};
\end{circuitikz}
\end{document}
